% Chapter 4
\chapter{Differential and integrated fiducial cross sections for Higgs boson production in the four-lepton decay channel}\label{Chapter4}
\chaptermark{Fiducial cross sections for $pp\rightarrow H \rightarrow  4\ell$}
\section{Introduction}
\label{sec:intro}

The ATLAS and CMS collaborations first reported the discovery of a new boson in 2012~\cite{Chatrchyan:2012ufa,Aad:2012tfa}
consistent with the standard model (SM) Higgs 
boson~\cite{Englert:1964et,Higgs:1964ia,Higgs:1964pj,Guralnik:1964eu,Higgs:1966ev,Kibble:1967sv} 
based on proton-proton ($\Pp\Pp$) collisions delivered by the CERN LHC at a center-of-mass energy of 
$\sqrt{s}=7\TeV$ in 2011 and $8\TeV$ in 2012. 
Subsequent studies by CMS using the full LHC Run 1 data set in various decay channels and production modes 
and combined measurements from ATLAS and CMS ~\cite{CMS:2014ega, AtlasProperties, Aad:2015zhl, CMS:2015kwa}
showed that the properties of the new boson are so far consistent with expectations for the SM Higgs boson.

The $\HZZfl$ decay channel ($\ell=\Pe,\Pgm$) has a large signal-to-background ratio due to the 
complete reconstruction of the final state decay products and excellent lepton momentum resolution
and is one of the most important channels for studies of the Higgs boson's properties. 
Measurements performed using this decay channel and the Run 1 data set include the determination of 
the mass and spin-parity of the new boson~\cite{CMSH4lLegacy,CMSH4lSpinParity,CMSH4lAnomalousCouplings}, 
its width~\cite{CMSH4lWidth,CMSH4lLifetime} and fiducial cross sections~\cite{CMSH4lFiducial8TeV}, as well as 
tests for anomalous HVV couplings~\cite{CMSH4lAnomalousCouplings,CMSH4lLifetime}. 

This analysis note presents measurements of properties of the Higgs boson in the $\HZZfl$ decay channel at $\sqrt{s}=13\TeV$
using $\usedLumiABC$ of $\Pp\Pp$ collision data collected with the CMS experiment at the LHC in 2016, 2017 and 2018.
The data collected in 2018 corresponding to $\usedLumiC$ is analyzed carefully and combined with $\usedLumiB$ of the data from 
2017~\cite{CMS-PAS-HIG-18-001} and with $\usedLumiA$ of data from 2016~\cite{CMS-HIG-16-041} that were previously analyzed. Therefore, the main
focus of this analysis note is to present the results of the analysis of 2018 data that is later on combined on a datacard level
with the previously analyzed 2016 and 2017 data.






\section{Datasets, simulation samples and theoretical predictions}\label{sec:DatasetsSimulation}
\subsection{Datasets}\label{sec:Datasets}
The dataset used in the measurements was collected by the CMS experiment in 2011 and 2012, and corresponds to \usedLumiA of 7~TeV collision data
and \usedLumiB of 8~TeV collision data respectively. All the details on the CMS standard selection of runs and luminosity sections, as well as the
selection of the triggers used to collect this dataset are described in Ref.~\cite{Chatrchyan:2013mxa}. Thus, similar detector operation conditions are 
imposed for the validation of the data to be used for the analysis of the $4e$, $4\mu$ and $2e2\mu$ channels. The absolute integrated luminosity of the
proton collisions in 2011 has been determined with a relative precision of 2.2\%~\cite{CMS:2011vpa}, while in 2012 it has been determined with a relative
precision of 2.6\%~\cite{CMS-PAS-LUM-13-001}. The analysis relies on primary datasets (PDs) produced centrally and which combine various collections of 
High Level Triggers (HLT).  The detailed content of the PDs evolves in phase with the evolution of the trigger menu to cope with ever increasing
instantaneous luminosity. The analysis relies on the so-called {\it DoubleElectron} and {\it DoubleMuon}. These latter PDs are formed by a {\tt OR} between
various triggers with symmetric or asymmetric trigger thresholds for the two leptons, with or without additional identification and isolation 
requirements. They also include triggers requiring three leptons above a low $p_T$ theshold.  It also includes cross-triggers which were added to
recover few percent of inefficiency in the 2e2$\mu$ channel at low Higgs boson masses, forming the so-called {\it MuEG} PD.  Tri-electron triggers for 2012
data are also used, in order to recover most of the inefficiency of DoubleElectron trigger.
All the triggers used with the 8~TeV data collected in 2012 are listed in Table~\ref{tab:triggers_2012}.
\begin{table}[h]
\scriptsize
    \centering
    \begin{tabular}{|l|l||l|l|c|}
\hline %-------------------------------------------------------------------------------------------------------------------------------------------
\# Channel     & Purpose & HLT path                                                                                & L1 seed            & prescale  \\ \hline
\hline %-------------------------------------------------------------------------------------------------------------------------------------------
        4e     & main    & \verb| HLT_Ele17_CaloTrk_Ele8_CaloTrk |                                 & \verb| L1_DoubleEG_13_7 |   & 1      \\
                     &            & \verb| OR HLT_Ele15_Ele8_Ele5_CaloIdL_TrkIdVL       |                      & \verb| L1_TripleEG_12_7_5 | & 1          \\
        4$\mu$ & main    & \verb| HLT_Mu17_Mu8  |                                                  & \verb| L1_Mu10_MuOpen   |   & 1       \\
               &         & \verb| OR HLT_Mu17_TkMu8   |                                            & \verb| L1_Mu10_MuOpen   |   & 1       \\
      2e2$\mu$ & main    & \verb| HLT_Ele17_CaloTrk_Ele8_CaloTrk   |                               & \verb| L1_DoubleEG_13_7 |   & 1      \\
               &         & \verb| OR HLT_Mu17_Mu8       |                                          & \verb| L1_Mu10_MuOpen   |   & 1      \\
               &         & \verb| OR HLT_Mu17_TkMu8  |                                             & \verb| L1_Mu10_MuOpen   |   & 1      \\
               &         & \verb| OR HLT_Mu8_Ele17_CaloTrk      |                                  & \verb| L1_MuOpen_EG12   |   & 1      \\
               &         & \verb| OR HLT_Mu17_Ele8_CaloTrk      |                                  & \verb| L1_Mu12_EG6      |   & 1      \\
\hline %-------------------------------------------------------------------------------------------------------------------------------------------
        4$\mu$ & backup  & \verb| HLT_TripleMu5    |                                               & \verb| L1_TripleMu0 |       & 1      \\
\hline %-------------------------------------------------------------------------------------------------------------------------------------------
4e and 2e2$\mu$ & Z T\&P         & \verb| HLT_Ele17_CaloTrkVT_Ele8_Mass50 |                        & \verb| L1_DoubleEG_13_7 |   & 5      \\
4e and 2e2$\mu$ & Z T\&P low pT  & \verb| HLT_Ele20_CaloTrkVT_SC4_Mass50_v1  |                     & \verb| L1_SingleIsoEG18er | & 10     \\
4$\mu$ and 2e2$\mu$ & Z T\&P     & \verb| HLT_IsoMu24_eta2p1|                                      & \verb| L1_SingleMu16er |    &        \\
4$\mu$ and 2e2$\mu$ & J/psi T\&P & \verb| HLT_Mu7_Track7_Jpsi   |                                  &                             &        \\
                    &            & \verb| HLT_Mu5_Track3p5_Jpsi  |                                 &                             &        \\
                    &            & \verb| HLT_Mu5_Track2_Jpsi   |                                  &                             &        \\
\hline %-------------------------------------------------------------------------------------------------------------------------------------------
    \end{tabular}
\small
    \caption{
    triggers used with the 8~TeV data collected in 2012.
    }
    \label{tab:triggers_2012}
\end{table}
\subsection{Simulation samples and theoretical predictions} \label{sec:Simulations}
Descriptions of the SM Higgs boson production in the gluon fusion ($\ggH$) process are obtained using 
the {\sc HRes}~\cite{deFlorian:2012mx, Grazzini:2013mca}, {\sc powheg}~\cite{powheg, powhegvv} and {\sc powheg minlo-HJ}~\cite{Hamilton:2012np} generators.
The {\sc HRes} generator is a partonic level Monte Carlo generator that provides a description of the $\ggH$ process at next-to-next-to-leading order 
(NNLO) accuracy in perturbative QCD and next-to-next-to-leading logarithmic (NNLL) accuracy in resummation of soft-gluon effects at small transverse 
momenta~\cite{deFlorian:2012mx, Grazzini:2013mca}. Since the resummation is inclusive over the QCD radiation recoiling against the Higgs
boson, {\sc HRes} is considered for the prediction of fiducial cross sections which are inclusive in the associated jet activity.
The {\sc powheg} generator is a partonic level matrix-element generator that implements next-to-leading (NLO) perturbative QCD calculations and 
additionally provides an interface to parton-shower programs. It provides a description of the $\ggH$ production in association with zero jets at NLO
accuracy. For the purpose of this analysis, it has been tuned to closely match the Higgs boson $p_{T}$ spectrum in the full phase space as predicted by
the {\sc HRes} generator. The {\sc powheg minlo-HJ} generator is an extension of the {\sc powheg} generator based on the MiNLO 
prescription~\cite{Hamilton:2012np} for the improved next-to-leading logarithmic (NLL) accuracy applied to the $\ggH$ production in association with up to
one additional jet. It provides a description of the $\ggH$ production in association with zero and one jets at the NLO accuracy, and the $\ggH$ production
in association with two jets only at the LO accuracy. Descriptions of the SM Higgs boson production in the vector boson fusion (VBF) process is obtained at
 NLO accuracy using also the {\sc powheg} generator.
The processes of SM Higgs boson production associated with gauge bosons (VH) or top-antitop quark pair ($\ttH$) are described at LO accuracy 
using {\sc pythia 6.4}~\cite{Sjostrand:2006za}, and higher order diagrams are accounted for by {\sc pythia}'s parton showering model. The Monte Carlo (MC) 
samples simulated with these generators are normalised to the inclusive SM Higgs boson production cross sections and branching fractions that correspond 
to the SM calculations at NNLO and NNLL accuracy, in accordance with the LHC Higgs Cross Section Working Group 
recommendations~\cite{LHCHiggsCrossSectionWorkingGroup:2013tqa}. In order to estimate the dependance of the measurement procedure on the underlying 
assumption for the Higgs boson production mechanism, we have used the set of MC samples for individual production mechanisms described in the paragraph 
above. In addition, in order to estimate the dependance of the measurement on different assumptions of the Higgs boson properties, we have also simulated
a range of samples that describe production and decay of exotic Higgs-like resonances to the four-lepton final state.
These include spin-zero, spin-one and spin-two resonances with anomalous interactions with a pair of neutral gauge bosons
($ZZ$, $Z\gamma^{*}$, $\gamma^{*}\gamma^{*}$) described by higher-order operators, as discussed in details in Ref.~\cite{Khachatryan:2014kca}. The 
procedure
used to estimate the model dependance effects will be presented in Section~\ref{sec:Systematics}. All these samples are generated using the {\sc powheg} 
generator for the description of NLO QCD effects in the production mechanism, and {\sc JHUgen}~\cite{Gao:2010qx,Bolognesi:2012mm,Anderson:2013afp} to 
describe the decay of these exotic resonances to four leptons including all spin correlations.
The list of Standard Model Higgs boson 8~TeV samples with corresponding values of cross sections can be found in Table~\ref{tab:MCsamplesSig}, while
all the signal Model Carlo samples studied in the analysis are listed in Table~\ref{tab:samplesSM}. A subset of these samples that is used to asses the
model dependance of the measurement results is indicated in the table.
In a similar way, the samples of the exotic spin-parity models and exotic production/decay modes are presented in Table~\ref{tab:samplesExo} in 
Appendix~\ref{app:sam}.
\begin{table}[!h!tb]
\begin{center}
\small
\caption{
Standard Model signal Model Carlo samples studies in the analysis.
Samples used for the estimation of the model dependance of the measurement results are denoted in the column ``D.''.
\label{tab:samplesSM}
}
\begin{tabular}{|l|c|c|c|} \hline
Sample & Description & $m_{H} (GeV)$ & D. \\ \hline
SMHiggsToZZTo4L\_M-125\_8TeV-powheg15-JHUgenV3-pythia6 & gg$\rightarrow$H ({\sc powheg+JHU}) & 125 $\GeV$ & Y \\
 SMHiggsToZZTo4L\_M-126\_8TeV-powheg15-JHUgenV3-pythia6 & gg$\rightarrow$H ({\sc powheg+JHU})& 126 $\GeV$ & Y \\
 GluGluToHToZZTo4L\_M-125\_8TeV-powheg15-pythia6 & gg$\rightarrow$H ({\sc powheg}) & 125 $\GeV$ & No \\
 GluGluToHToZZTo4L\_M-126\_8TeV-powheg15-pythia6 & gg$\rightarrow$H ({\sc powheg}) & 126 $\GeV$ & No \\
 GluGluToHToZZTo4L\_M-125\_8TeV-powheg-pythia6 & gg$\rightarrow$H ({\sc powheg1.0}) & 125 $\GeV$ & No \\
 GluGluToHToZZTo4L\_M-126\_8TeV-powheg-pythia6 & gg$\rightarrow$H ({\sc powheg1.0}) & 126 $\GeV$ & No \\
 GluGluToHToZZTo4L\_M-125\_8TeV-minloHJJ-pythia6-tauola & gg$\rightarrow$H ({\sc minloHJJ}) & 125 $\GeV$ & Y\\
 GluGluToHToZZTo4L\_M-126\_8TeV-minloHJJ-pythia6-tauola & gg$\rightarrow$H ({\sc minloHJJ}) & 126 $\GeV$ & Y\\
 VBF\_HToZZTo4L\_M-125\_8TeV-powheg-pythia6 & VBF ({\sc powheg}) & 125 $\GeV$ & Y \\
 VBF\_HToZZTo4L\_M-126\_8TeV-powheg-pythia6 & VBF ({\sc powheg}) & 126 $\GeV$ & Y \\
 WH\_HToZZTo4L\_M-125\_8TeV-pythia6 & WH ({\sc pythia}) & 125 $\GeV$ & Y \\
 WH\_HToZZTo4L\_M-126\_8TeV-pythia6 & WH ({\sc pythia}) & 126 $\GeV$ & Y \\
 ZH\_HToZZTo4L\_M-125\_8TeV-pythia6 & ZH ({\sc pythia}) & 125 $\GeV$ & Y \\
 ZH\_HToZZTo4L\_M-126\_8TeV-pythia6 & ZH ({\sc pythia}) & 126 $\GeV$ & Y \\
 TTbarH\_HToZZTo4L\_M-125\_8TeV-pythia6 & ttH ({\sc pythia})& 125 $\GeV$ & Y \\
 TTbarH\_HToZZTo4L\_M-126\_8TeV-pythia6 & ttH ({\sc pythia})& 126 $\GeV$ & Y \\
\hline
\end{tabular}
\normalsize
\end{center}
\end{table}
The MC event samples that are used to estimate contribution from the background process $\Pg\Pg \rightarrow \cPZ\cPZ$ were simulated using
{\sc MCFM 6.7}, while the background process $\Pq\Paq \to 4\ell$ was simulated at NLO accuracy with {\sc powheg} generator including s-, t- and u-channel
diagrams. For the purpose of the $Z\rightarrow4\ell$ cross-section measurements, we also have separately modelled contributions from the t- and u-channels
of the $\Pq\Paq(\to \Z\Z*)\to 4\ell$ process at NLO accuracy with {\sc powheg} generator.
All event generators described above take into account the internal initial-state and final-state radiation effects which can lead to the
presence of additional hard photons in an event.
The {\sc powheg} and {\sc JHUgen} event generators also properly take into account interference between all contributing amplitudes,
 including those related to the permutations of identical leptons in the $4e$ and $4\mu$ final states.
In case of the LO generators, the set of parton distribution functions (PDF) \textsc{cteq6l}~\cite{cteq66} is used,
in case of the NLO generators the PDF set \textsc{CT10}~\cite{CT10} is used,
while the PDF set \textsc{MSTW2008}~\cite{MSTW08} is used in case of the NNLO generator.
All generated events are interfaced to {\sc pythia 6.4} Tune $Z2^{*}$~\cite{Chatrchyan:2011id} to simulate the effects of the parton shower, hadronization
and underlying event.
The {\sc pythia 6.4} Z2* tune is derived from the Z1 tune~\cite{Field:2010bc}, which uses the CTEQ5L parton distribution set, whereas Z2* adopts 
CTEQ6L~\cite{Pumplin:2002vw}.
In case of {\sc HRes}, since it does not provide an interface to programs that can simulate the hadronization and underlying event effects, the size of
these effects is first estimated using the {\sc powheg+JHUgen} and its interface to {\sc Pythia 6.4}, and then the correction is applied to the {\sc HRes}
parton level predictions.
The generated events were processed through a detailed simulation of the CMS detector based on
\GEANTfour~\cite{GEANT} and were reconstructed with the same algorithms
that are used for data analysis.
Multiple $\Pp\Pp$ interactions (pileup) are included in simulations to match the distribution of the number of
interactions per LHC bunch crossing observed in data.  The average
number of pileup interactions is estimated to be approximately 9 and 21 in the 7 and 8 \TeV datasets, respectively.
Energy deposits from pileup interactions and from the underlying event are subtracted
when computing the energy of jets and isolation of reconstructed objects
using the \textsc{FastJet} technique~\cite{Cacciari:2007fd,Cacciari:2008gn,Cacciari:2011ma}.
This technique is based on the calculation of the $\eta$-dependent average energy density %($\rho$), 
evaluated individually for each event.

The selection efficiency in all simulated samples is rescaled in order to correct for residual differences in lepton selection efficiencies in data and 
simulation. This correction is based on the total lepton selection efficiencies measured in inclusive samples of Z-boson events in simulation and data 
using a ``tag and probe'' method~\cite{CMS:2011lsa}, separately for 7 and 8~TeV collisions.
%The deviation of the efficiency in simulation relative to data, for the majority of the phase space of the leptons, is less than 3\% for both electrons and muons. 
The dependency of the selection efficiency on the number of reconstructed primary vertices in the event is negligible for both the 7 and 8 TeV data samples.
More details are presented and discussed in Ref.~\cite{Chatrchyan:2013mxa}.


%%%%%%%%%
%
\begin{table}[htb]
%
  \begin{center}
  \vspace*{0.3cm}

    \begin{tabular}{llll}
\hline %---------------------------------------------------------------------------------------------------------------------------------------------------------------------------- 
\small      \bf{Process}                & \small \bf{MC}                        & \small\boldmath{$\sigma_{(N)NLO} $} &  Comment \\
\hline \hline%==============================================================================================
\small          gg$\rightarrow {\rm H}(\rightarrow {\rm Z}{\rm Z}) \rightarrow 4\ell$                   & \small{\sc powheg+JHUgen}     & \small 19.27  pb      & gg$\rightarrow$H production\\
\small          ${qq} \rightarrow {\rm H}(\rightarrow {\rm Z}{\rm Z})+qq \rightarrow 4\ell+$X   & \small{\sc powheg}    & \small 1.578 pb       & VBF production\\
\small          ${qq} \rightarrow {\rm H}(\rightarrow {\rm Z}{\rm Z})+W \rightarrow 4\ell+$X    & \small{\sc pythia}    & \small  0.7046 pb     & WH production\\
\small          ${qq} \rightarrow {\rm H}(\rightarrow {\rm Z}{\rm Z})+Z \rightarrow 4\ell+$X            & \small{\sc pythia}    & \small 0.4153 pb      & ZH production\\
\small          ${qq} \rightarrow {\rm H}(\rightarrow {\rm Z}{\rm Z})+t\bar{t} \rightarrow 4\ell+$X     & \small{\sc pythia}    & \small 0.1293 pb      & ttH production\\
%\hline %---------------------------------------------------------------------------------------------------------------------------------------------------------------------------- 
%\small         ${gg} \rightarrow {\rm H} \rightarrow {\rm Z}{\rm \gamma*} \rightarrow 4\ell $          & \small{\tt powheg + JHUGen}   & \small ?  fb  & Z$\gamma*$ decay\\
%\small         ${gg} \rightarrow {\rm H} \rightarrow {\rm \gamma*}{\rm \gamma*} \rightarrow 4\ell $    & \small{\tt powheg + JHUGen}   & \small ?  fb  & $\gamma*\gamma*$ decay\\
\hline %---------------------------------------------------------------------------------------------------------------------------------------------------------------------------- 
   \end{tabular}

  \end{center}
%
  \small
  \caption{  MC simulation datasets used for the signal processes predicted by the SM for m(H)=125.0$\GeV$ (${\rm Z}$ stands for ${\rm Z}$, ${\rm Z}^*$, $\gamma^*$). The cross sections do not include the ZZ and $4\ell$ branching ratios. In the SM, the branching ratio ${\rm H} \rightarrow {\rm Z}{\rm Z}\rightarrow 4\ell$ is 1.247$\times 10^{-4}$ for m(H)=125.0 $\GeV$.}
  \label{tab:MCsamplesSig}
%
\end{table}
%%%%%%
\begin{table}[htb]
%
  \begin{center}
  \vspace*{0.3cm}

    \begin{tabular}{llll}
\hline %---------------------------------------------------------------------------------------------------------------------------------------------------------------------------- 
\small      \bf{Process}                & \small \bf{MC}                        & \small\boldmath{$\sigma_{(N)NLO} $} &  Comment \\
\hline \hline%==============================================================================================
\small          $ q\bar{q} \rightarrow {\rm Z}{\rm Z} \rightarrow 4e(4\mu,4\tau) $        & \small{\sc powheg}  & \small 76.91 fb       & s+t/u channels\\
\small  $ q\bar{q} \rightarrow {\rm Z}{\rm Z} \rightarrow 2e2\mu  $             & \small{\sc powheg}            & \small 176.7 fb       & s+t/u channels\\
\small  $ q\bar{q} \rightarrow {\rm Z}{\rm Z} \rightarrow 2e(2\mu)2\tau  $       & \small{\sc powheg}           & \small 176.7 fb       & s+t/u channels\\
\small          $ q\bar{q} \rightarrow {\rm Z}{\rm Z} \rightarrow 4e $                           & \small{\sc powheg}   & \small 494.8 fb       & t/u channels only, $m_{ll}>1GeV$\\
\small  $ q\bar{q} \rightarrow {\rm Z}{\rm Z} \rightarrow 2e2\mu  $             & \small{\sc powheg}            & \small 1527 fb        & t/u channels only, $m_{ll}>1GeV$\\
\small  $ q\bar{q} \rightarrow {\rm Z}{\rm Z} \rightarrow 4\mu  $                & \small{\sc powheg}   & \small 488.4 fb       & t/u channels only, $m_{ll}>1GeV$\\
\small          ${gg} \rightarrow {\rm Z}{\rm Z} \rightarrow 2e2\mu $                           & \small{\sc MCFM 6.7}          & \small 2.444  fb      & K-factor of 2.66 included\\
\small          ${gg} \rightarrow {\rm Z}{\rm Z} \rightarrow 4\mu$                              & \small{\sc MCFM 6.7}          & \small 1.222 fb       & K-factor of 2.66 included\\
\small          ${gg} \rightarrow {\rm Z}{\rm Z} \rightarrow 4e$                                        & \small{\sc MCFM 6.7}          & \small 1.222 fb       & K-factor of 2.66 included\\
\hline %---------------------------------------------------------------------------------------------------------------------------------------------------------------------------- 
   \end{tabular}

  \end{center}
%
  \small
  \caption{  MC simulation datasets used for the background processes (${\rm Z}$ stands for ${\rm Z}$, ${\rm Z}^*$, $\gamma^*$);
  In MC samples with label ``t/u channels only'' only the t/u channels are simulated in the $qq \to 4\ell$ process and these samples are used for the $Z \to 4\ell$ analysis.
  The ``t/u channels only'' samples have lower cut in $m_{ll}$ of 1 GeV, while all other $qq \to ZZ$ and $gg \to ZZ$ samples have lower cut in $m_{ll}$ of 4 GeV. In case of the $gg \to ZZ$ background, the overall NLO k-factor of 2.66 is included, as discussed in Section~\ref{sec:EventsBackgrounds}.}
  \label{tab:MCsamplesBkg}
%
\end{table}

\section{Analysis strategy}\label{sec:AnalysisStrategy} 
\subsection{Event selection and background models} 
\subsection{Definition of the fiducial volume}
\subsection{Signal model and statistical procedure}
\subsection{Extraction of inclusive \texorpdfstring{$\Hllll$}{H to 4l} fiducial cross section}
\subsection{Extraction of differential \texorpdfstring{$\Hllll$}{H to 4l} fiducial cross sections}
\section{Systematic uncertainties}\label{sec:Systematics}
\section{Results}\label{sec:Results}
\subsection{Measurement of inclusive fiducial cross section}
\subsection{Measurement of differential fiducial cross sections}
\section{Summary}\label{sec:Summary}

%===========================================================================
